\subsection{Tasks, Routes, and Motion}
Let $\mathcal V=\{1,\ldots,N\}$ contain task locations $z_i\in\mathbb R^2$ and nonnegative rewards $r_i$. Agents $m\in\mathcal A=\{1,\ldots,M\}$ share start $z_S$ and destination $z_E$. Agent $m$ has a flight-time budget $B_m$, initial velocity $v_m^S$, and componentwise velocity and acceleration bounds $\bar v_{md}$ and $\bar a_{md}$ for $d\in\{x,y\}$. We ignore collisions, obstacles, and jerk limits.

An agent selects an ordered sequence of \emph{visit occurrences} $\pi_m=(\nu_{m0},\ldots,\nu_{m,K_m-1})$ beginning at START and ending at END. An interior occurrence is TASK, with fixed position $z_i$ and an associated task ID, or FREE, with a chosen position and no intrinsic reward. Occurrences have distinct identities even when positions or task IDs coincide. Each has a boundary velocity $v_{mk}$ and a nonnegative pre-departure wait $w_{mk}$. End velocity is zero and a positive wait requires zero velocity. An explicit TASK occurrence earns reward even if traversed at nonzero velocity; incidental flyovers and FREE visits at a task position do not.

The motion joining consecutive visits obeys
\begin{equation}
 \dot p_m(t)=v_m(t),\quad \dot v_m(t)=a_m(t),\quad
 |v_{md}(t)|\le\bar v_{md},\quad |a_{md}(t)|\le\bar a_{md}.
 \label{eq:dynamics}
\end{equation}
The implemented connector represents each connection by finitely many constant-acceleration pieces and chooses a minimum-time feasible connection for fixed endpoint states. If $D_{mk}$ is its duration, arrival and departure times satisfy
\begin{equation}
 \alpha_{m0}=0,\quad \beta_{mk}=\alpha_{mk}+w_{mk},\quad
 \alpha_{m,k+1}=\beta_{mk}+D_{mk}.
 \label{eq:timing}
\end{equation}
An agent's expenditure and the mission horizon are
\begin{equation}
 c_m=\alpha_{m,K_m-1}\le B_m,\qquad
 \bar c=\frac1M\sum_m c_m,\qquad T=\max_m c_m.
 \label{eq:budget}
\end{equation}
Waiting before END consumes budget. After arrival the agent remains at $z_E$ until $T$, contributing to communication without further expenditure. Flight time is a proxy for resource use, not a physical energy model.

Let $\mathcal Q_m$ be the task IDs in agent $m$'s explicit TASK occurrences. Repeated or shared visits earn each reward once:
\begin{equation}
 Q(S)=\sum_{i\in\cup_m\mathcal Q_m} r_i.
 \label{eq:reward}
\end{equation}

\subsection{Range Contacts and Information Propagation}
At time $t$, a pair $(m,n)$ is within range if $\|p_m(t)-p_n(t)\|_2\le R$. For each maximal closed within-range interval $[s,e]$, the usable link interval is $[s+\delta,e]$ when $s>0$ and $e-s\ge\delta$, where $\delta$ is the continuous handshake duration. The initial interval beginning at zero is usable immediately because initial handshakes and synchronization are assumed complete. A later interruption resets the handshake. Thus communication can occur during movement, waits, or post-arrival presence; a planned meeting is sufficient but not necessary.

Agents retain and forward information. Simultaneously linked agents exchange the latest data across each connected component, while links at different times must respect chronological order. Let $h_{nm}(t)$ be the latest timestamp of information originating at source $m$ available to receiver $n$ by $t$, with $h_{nm}(0)=0$ and $h_{mm}(t)=t$. The directed information age and its mission peak are
\begin{equation}
 A_{nm}(t)=t-h_{nm}(t),\qquad
 \Delta(S)=\sup_{0\le t\le T}\max_{m\ne n} A_{nm}(t).
 \label{eq:staleness}
\end{equation}
The supremum includes left limits immediately before a contact, when age may be larger than just after information exchange. For one agent we set $\Delta=0$. The communication requirement is $\Delta(S)\le T_c$. Equivalently, under initial synchronization, every full mission window of length $T_c$ must permit time-respecting propagation of every source's information to every other agent. If $T<T_c$, no full window is required. No additional delivery condition is imposed after $T$.

\subsection{Objectives and Feasibility}
We minimize the three-component vector
\begin{equation}
 f_{T_c}(S)=\left(-Q(S),\;\bar c(S),\;V_c(S)\right),\quad
 V_c(S)=\max\{0,\Delta(S)-T_c\}.
 \label{eq:objectives}
\end{equation}
The flight-time bounds in~\eqref{eq:budget}, dynamics in~\eqref{eq:dynamics}, and visit rules are hard physical constraints. A physically valid candidate with $V_c>0$ may be retained during search as an infeasible communication stepping stone. The reported feasible reward--time front is the $V_c=0$ slice. All candidates below $T_c$ tie on the third objective, so the final front does not reward unnecessary communication margin.
